\documentclass[10pt,a4paper]{amsart}
\usepackage[margin=19mm]{geometry}
\usepackage{amsmath,amssymb,mathtools}
\usepackage[scr=rsfs,cal=euler]{mathalfa}
\usepackage{xfrac}
\usepackage[unicode,hidelinks,bookmarks=false]{hyperref}
\newcommand{\ie}{i.e.\ }
\newcommand{\cf}{cf.\ }
\newcommand{\resp}{respectively\ }
\newcommand{\Subsection}{\subsection}
\providecommand{\nobreakdash}{\nobreak\hbox{-}}
\setlength{\emergencystretch}{2em}
\multlinegap0pt
\usepackage{bm}
\usepackage{bbm}
\usepackage{dsfont}
\usepackage{xspace}
\usepackage{cancel}
\usepackage{mathtools}
\usepackage{amsthm}
\makeatletter
\@ifundefined{theorem}{\newtheorem{theorem}{Theorem}}{}
\@ifundefined{cor}{\newtheorem{cor}[theorem]{Corollary}}{}
\makeatother
\DeclareMathOperator{\fin}{FIN}
\newcommand{\UU}{\mathcal{U}}
\title[Infinite random graphs]{Infinite random graphs}
\author{Ziemowit Kostana, Jaros{\l}aw Swaczyna, Agnieszka Widz}
\date{Source: arXiv:2601.16013v1 (CC BY 4.0)}
\begin{document}
\maketitle

\noindent\emph{Source and licence.} Infinite random graphs. Version arXiv:2601.16013v1 (CC BY 4.0), distributed under the Creative Commons Attribution 4.0 International licence, as recorded in the arXiv licence metadata for this exact version and shown on the abstract page https://arxiv.org/abs/2601.16013v1. Full rights notice: licenses/arxiv-2601.16013-CC-BY-4.0.txt.\par\smallskip

\noindent\textit{Editorial scope.} Counted unit: the Theorem whose source label is \texttt{t:BasisThm} in the article (source0.tex lines 1586--1586, source label \texttt{t:BasisThm}), delivered together with the article's own complete original proof of it (source0.tex lines 1587--1672, one \texttt{proof} environment of 86 lines). No mathematics is written, repaired or re-derived: statement and proof are the article's own text line for line. Required context retained uncounted: partitionregular (lines 1569--1573). Every definition, display and label of those lines is retained unchanged. Omitted from this package and not redistributed: 1--1568, 1574--1585, 1673--1843. Nothing inside a retained excerpt is omitted or altered: every retained body is delivered exactly as the article's lines are written, so no omission receipt of any kind accompanies this package. Citations: the retained excerpts carry no citation command at all, so no bibliography entry of the article is transcribed and no reference is redistributed. Reference adaptation: none was needed and none was made; every reference inside the delivered unit resolves inside it, so no unresolved reference or citation is produced. Printed slips kept literal: no printed word, symbol, hypothesis or number of the counted statement or of its proof was altered. Layout and class: the article's own class is \texttt{amsart} with packages including amssymb, amsthm, mathtools, babel, fontenc, inputenc, color, graphicx, bbm, xcolor, physics, amsmath, tikz, wrapfig; the wrapper is the lane's pinned \texttt{amsart} editorial wrapper at 10pt on A4, which restates the theorem-like environments the retained text needs and adds the article's own macro definitions that the retained text needs (\texttt{\textbackslash UU}, \texttt{\textbackslash fin}), with extra wrapper packages (bm, bbm, dsfont, xspace, cancel, mathtools). The delivered numbering is the unit's own. Assets: no retained span contains a figure environment, a graphics-inclusion command, a PDF-page inclusion, a table or a TikZ picture; no image file is copied into this package, the package declares no asset dependency and no image file is redistributed with it. Licence: version arXiv:2601.16013v1 of this article is distributed under the Creative Commons Attribution 4.0 International licence, as recorded in the arXiv licence metadata for this exact version and shown on the abstract page https://arxiv.org/abs/2601.16013v1. A full rights notice is delivered at licenses/arxiv-2601.16013-CC-BY-4.0.txt. Characters: every retained excerpt is pure ASCII, so the delivered engine, XeLaTeX with its default Latin Modern text font, reports no missing character.
\begin{cor}\label{partitionregular}
    If $(\omega,\,E)$ is any weakly universal graph, then for any partition
    $$\omega=A_0\cup \ldots \cup A_{k-1}$$
    there is some $i<k$ such that the induced subgraph $(\omega, \, E)\restriction_{A_i}$ is weakly universal.
\end{cor}

\begin{theorem}[Basis Theorem]\label{t:BasisThm} If $G$ is weakly universal, then $G$ contains either $\UU_{\fin}$ or its complement. \end{theorem}

\begin{proof}
Let \((\omega, E)\) be weakly universal.  
By iteratively applying Corollary~\ref{partitionregular}, we recursively define a function
\(f\colon \omega \longrightarrow \{0,1\}\) so that, for each \(n<\omega\), the set


    \[
U_n := 
\left\{\,k\ge n :
  \mathrel{\vcenter{\hbox{%
    $\displaystyle
      \mathop{\scalebox{1.1}{$\forall$}}_{\,\scriptstyle i<n}$
  }}}%
  \!\!\!\;\bigl(\{i,k\}\in E \iff f(i)=1\bigr)
\right\}.
\]

is weakly universal.  
The construction is straightforward.

Applying Corollary \ref{partitionregular}, we see that for each $n<\omega$, at least one of the sets


\[
U_{n}\cap f^{-1}(\{0\})
\quad\text{and}\quad
U_{n}\cap f^{-1}(\{1\})
\]


is weakly universal.  
Evidently, one of these alternatives occurs infinitely often; and since the sequence \((U_{n})_{n<\omega}\) is decreasing, it in fact occurs for every \(n<\omega\).  
Without loss of generality, assume that all sets of the form
\[
U_{n}\cap f^{-1}(\{0\})
\]

are weakly universal.

Let \(\{H_{n} : n<\omega\}\) be a list of all finite graphs, up to isomorphism.  
We recursively choose finite sets \(K_{m}\subseteq \omega\) satisfying for all \(m<\omega\):
\begin{enumerate}
  \item \(K_{m}\subseteq f^{-1}(\{0\})\),
  \item \((\omega, E)\restriction_{K_{m}} \simeq H_{m}\),
  \item $K_m$ is disjoint and disconnected from $\bigcup\limits_{i<m}K_i.$
\end{enumerate}

Suppose the sets \(K_{i}\) are chosen for all \(i<m\).  
Let \(n<\omega\) be large enough so that \(\bigcup_{i<m} K_{i} \subseteq n\).  
Since \(U_{n}\cap f^{-1}[\{0\}]\) is weakly universal, we may choose


\[
K_{m} \subseteq U_{n}\cap f^{-1}(\{0\})
\quad\text{so that}\quad
(\omega, E)\restriction_{K_{m}} \simeq H_{m}.
\]


To verify (3), observe that (1) together with the definition of \(U_{n}\) ensures


\[
\Bigl\{\,k \in \omega\setminus n : k \text{ is disconnected from } \bigcup_{i<m} K_{i}\Bigr\}
\;\supseteq\; U_{n}.
\]


This completes the recursive construction.  
By (2) and (3) we obtain


$$(\omega,\, E)\restriction_{\bigcup\limits_{m<\omega}K_m} \simeq \UU_{\fin}.$$



The other, symmetric case occurs when the sets


\[
U_{n}\cap f^{-1}(\{1\})
\]


are weakly universal.  Then the analogous construction gives us a complement of $\UU_{\fin}$.
\end{proof}

\end{document}
